% Two state machines side by side, and the emergency frame's layout.
\begin{tikzpicture}[font=\sffamily\scriptsize, x=1cm, y=1cm]

% ================================================================ the node's states
\node[chlabel, anchor=west] at (-6.6,3.6) {the node's own state: five, and small on purpose};
\node[initial] (i) at (-6.4,2.4) {};
\node[runstate, text width=17mm] (boot) at (-5.2,2.4) {BOOT};
\node[runstate, text width=17mm] (pre)  at (-2.6,2.4) {PRE-OP};
\node[runstate, text width=17mm] (op)   at (0.0,2.4) {OPERATIONAL};
\node[degraded, text width=17mm] (deg)  at (0.0,0.6) {DEGRADED};
\node[safestate, text width=17mm] (stop) at (-2.6,0.6) {STOPPED};

\draw[trans] (i) -- (boot);
\draw[trans] (boot) -- node[above]{self-check} (pre);
\draw[trans] (pre) -- node[above]{fresh command} (op);
\draw[trans] (op) -- node[right]{any one of three} (deg);
\draw[trans] (deg) -- node[below]{bus-off} (stop);
\draw[recoveredge] (stop) to[bend left=30] node[below]{recovery, by policy} (pre);

\node[tlabel, anchor=north west, text width=44mm] at (1.4,1.4)
  {The three routes out of operational are independent: the master is absent,
   the command expired, or the controller went error-passive. Any one is
   enough.};

% ================================================================ the controller's states
\node[chlabel, anchor=west] at (-6.6,-0.9) {and the controller's own, which change with no software at all};
\foreach \i/\n/\sty in {0/{ERROR-ACTIVE}/runstate, 1/{ERROR-WARNING}/degraded,
                        2/{ERROR-PASSIVE}/degraded, 3/{BUS-OFF}/safestate}{
  \node[\sty, text width=20mm] (c\i) at ({-5.4+\i*2.9},-1.9) {\n};
}
\foreach \i in {0,1,2}{ \draw[trans] (c\i) -- (c\the\numexpr\i+1\relax); }
\draw[recoveredge] (c3) to[bend right=25] node[above]{only by an explicit restart} (c0);
\node[tlabel, anchor=north, text width=100mm, align=center] at (-1.0,-2.6)
  {counters rise on errors and fall on successful frames. The thresholds are in
   the data link standard, which this volume did not fetch, so the node reports
   the counters themselves rather than printing values it has not read};

% ================================================================ the emergency frame
\begin{scope}[yshift=-52mm, xshift=-64mm]
  \node[signame] at (-2mm,3mm) {emergency};
  \framefield{0}{24}{fid}{error code\\{\tiny 2 bytes}}
  \framefield{24}{16}{fcrc}{error register\\{\tiny 1}}
  \framefield{40}{16}{fstuff}{which condition\\{\tiny 1}}
  \framefield{56}{40}{fdata}{additional information, 4 bytes}
  \node[chlabel, anchor=west] at (100mm,3mm) {the profile's layout, worth following};
\end{scope}

\node[chlabel, anchor=north west, text width=54mm] at (-6.6,-6.4)
  {The code ranges, from the profile: 0x10xx generic, 0x20xx current, 0x30xx
   voltage, 0x40xx temperature, 0x50xx hardware, 0x60xx software, 0x80xx
   monitoring, 0xFFxx device specific. Following them costs nothing and means
   another vendor's tool can read this node's faults.};
\node[umlnote, text width=54mm, anchor=north west] at (0.4,-6.4)
  {This frame is sent once per event. The fault word in chapter 11's state
   frame is the other shape, sent every period, so a listener that arrives late
   still learns what is true. A joint node wants both, and the two are not
   redundant: one is a record of an event, the other is a statement of a
   condition.};
\end{tikzpicture}
